% !TEX program = lualatex
\documentclass[lualatex,ja=standard]{bxjsarticle}

\usepackage{amsmath}   % 高度な数式環境
\usepackage{amssymb}   % 数学記号

\title{数式のサンプル}
\author{LaTeX サンプル}
\date{\today}

\begin{document}

\maketitle

% -------------------------------------------------------------------
% インライン数式
% -------------------------------------------------------------------
\section{インライン数式}

文中に数式を入れる場合は \( E = mc^2 \) のように \texttt{\textbackslash(...\textbackslash)} を使います。
ドル記号 \$ でも同じ結果になります: $a^2 + b^2 = c^2$

% -------------------------------------------------------------------
% ディスプレイ数式
% -------------------------------------------------------------------
\section{ディスプレイ数式}

独立した行に大きく表示するには \texttt{\textbackslash[...\textbackslash]} を使います:

\[
  \int_{-\infty}^{\infty} e^{-x^2} \, dx = \sqrt{\pi}
\]

番号付きにするには \texttt{equation} 環境を使います:

\begin{equation}
  \sum_{n=1}^{\infty} \frac{1}{n^2} = \frac{\pi^2}{6}
  \label{eq:basel}
\end{equation}

式~\eqref{eq:basel} はバーゼル問題の解です。

% -------------------------------------------------------------------
% 分数・指数・根号
% -------------------------------------------------------------------
\section{分数・指数・根号}

分数: $\dfrac{a+b}{c-d}$

指数: $x^{2n}$、下付き: $a_{ij}$、組み合わせ: $x_i^2$

平方根: $\sqrt{2}$、n乗根: $\sqrt[n]{x}$

% -------------------------------------------------------------------
% ギリシャ文字
% -------------------------------------------------------------------
\section{ギリシャ文字}

小文字: $\alpha, \beta, \gamma, \delta, \epsilon, \theta, \lambda, \mu, \pi, \sigma, \phi, \omega$

大文字: $\Gamma, \Delta, \Theta, \Lambda, \Pi, \Sigma, \Phi, \Omega$

% -------------------------------------------------------------------
% 数学演算子・記号
% -------------------------------------------------------------------
\section{数学記号}

関係: $\leq,\ \geq,\ \neq,\ \approx,\ \equiv,\ \propto$

集合: $\in,\ \notin,\ \subset,\ \subseteq,\ \cup,\ \cap,\ \emptyset$

矢印: $\to,\ \Rightarrow,\ \Leftrightarrow,\ \infty$

% -------------------------------------------------------------------
% 複数行の数式（align）
% -------------------------------------------------------------------
\section{複数行の数式}

\texttt{align} 環境で等号を揃えて整列できます:

\begin{align}
  (a + b)^2 &= a^2 + 2ab + b^2 \\
  (a - b)^2 &= a^2 - 2ab + b^2 \\
  (a + b)(a - b) &= a^2 - b^2
\end{align}

% -------------------------------------------------------------------
% 行列
% -------------------------------------------------------------------
\section{行列}

\[
  A = \begin{pmatrix}
    a_{11} & a_{12} & a_{13} \\
    a_{21} & a_{22} & a_{23} \\
    a_{31} & a_{32} & a_{33}
  \end{pmatrix}
\]

各括弧の種類: \texttt{pmatrix}（丸）、\texttt{bmatrix}（角）、\texttt{vmatrix}（縦棒）

\[
  \det(A) = \begin{vmatrix}
    a & b \\
    c & d
  \end{vmatrix} = ad - bc
\]

% -------------------------------------------------------------------
% 場合分け
% -------------------------------------------------------------------
\section{場合分け}

\[
  |x| = \begin{cases}
    x  & (x \geq 0) \\
    -x & (x < 0)
  \end{cases}
\]

\end{document}